\chapter*{Tóm tắt}
\label{summary}

Các mô hình Encoder-only như BERT giữ vai trò nền tảng cho các bài toán hiểu ngôn ngữ
tự nhiên, nhưng trong nhiều năm gần như không được cập nhật về kiến trúc. NeoBERT
(2025) thu hẹp khoảng cách này bằng cách tích hợp các cải tiến từ thế hệ mô hình ngôn
ngữ lớn -- RoPE, SwiGLU, Pre-RMSNorm, tỷ lệ sâu--hẹp tối ưu và ngữ cảnh 4.096 token --
song chỉ được tiền huấn luyện trên tiếng Anh, trong khi các Encoder tiếng Việt hiện có
vẫn dựa trên kiến trúc thế hệ cũ.

Khóa luận xây dựng \textbf{ViNeoBERT} -- mô hình NeoBERT cho tiếng Việt -- bằng
một quy trình chuyển đổi ngôn ngữ thay vì huấn luyện lại từ đầu: (1)~khởi tạo lớp véc-tơ đầu vào và lớp đầu ra bằng kỹ thuật Chuyển đổi tuyến tính trong không gian có ngữ nghĩa (Semantic Aware Linear Transfer -- SALT), tái sử dụng
không gian véc-tơ của mô hình tiếng Việt ViDeBERTa thông qua các ánh xạ tuyến tính cục
bộ theo từng token; (2)~căn chỉnh với phần thân đóng băng: đóng băng toàn bộ phần thân Transformer,
chỉ huấn luyện hai ma trận vừa khởi tạo cho khớp với phần thân đó; và (3)~tiếp tục tiền huấn luyện trên ngữ liệu CulturaX tiếng
Việt theo lịch điều chỉnh tốc độ học Warmup--Stable--Decay. So với SALT gốc, đề tài đưa ra ba cải
tiến: mở rộng tập điểm neo bằng các cặp dịch Anh--Việt; khởi tạo riêng lớp đầu ra --
khởi tạo hệ số điều chỉnh theo tần suất xuất hiện của token trong tiếng Việt, đồng thời thu nhỏ
trọng số lớp đầu ra để tín hiệu tần suất này chi phối giai đoạn huấn luyện đầu; và bổ sung
giai đoạn căn chỉnh với phần thân đóng băng.

Thực nghiệm so sánh có kiểm soát trên cùng một lượng token huấn luyện xác nhận từng lựa chọn
thiết kế, từ việc giảm biên độ trọng số lớp đầu ra tới khởi tạo riêng hai ma trận và giai đoạn căn chỉnh với phần thân đóng băng. Trên bộ đánh giá gồm mười hai tác vụ hiểu tiếng Việt (phân loại, suy
luận, gán nhãn chuỗi, tương đồng ngữ nghĩa và biểu diễn câu), ViNeoBERT sau 5 tỷ token
($\sim$20GB) đạt kết quả tương đương PhoBERT-base-v2 ở đa số tác vụ và dẫn đầu ở gán nhãn
từ loại. Trên tác vụ đọc hiểu với bộ dữ liệu UIT-ViQuAD, mô hình đạt
76,6 điểm F1 -- hơn XLM-R-base 3,3 điểm và cạnh tranh với PhoBERT-large (lớn gấp 1,5 lần)
dù chỉ dùng khoảng 3,5\% lượng token mà mô hình này đã huấn luyện ($\sim$140 tỷ).
Ngoài ra, một pha mở rộng ngắn $\sim$200 triệu token ở độ dài 4.096 token giữ
pseudo-perplexity ổn định trên toàn dải ngữ cảnh, biến trần 4.096 token của kiến trúc
thành năng lực sử dụng được. Kết quả cho thấy tái sử dụng không gian biểu diễn sẵn có
là con đường tiết kiệm và hiệu quả để đưa các kiến trúc Encoder thế hệ mới đến với
tiếng Việt.

\vspace{4mm}
\noindent\textbf{Từ khóa:} NeoBERT, tiếng Việt, chuyển đổi ngôn ngữ, khởi tạo véc-tơ
cho ngôn ngữ mới, SALT, tiếp tục tiền huấn luyện.
